\chapter{Learning from repeated experiments}
\label{ch:repetitions}

\section{One history cannot tell us a probability}

A single run of the model produces one history. The focal allele might disappear. It might eventually occupy the whole population. Watching that one history does not tell us how often either outcome occurs.

Think about tossing a coin. Seeing one head does not establish that every future toss will be a head. To learn about the chance of a head, we repeat the experiment. We use the same idea here, except that one experiment is an entire population history, possibly containing thousands of reproduction events.

A \emph{replicate} means one fresh repetition of the whole experiment. Every replicate starts from the same population size, initial allele count and fitness advantage. It uses a fresh part of the random-number sequence. It continues until the allele is lost or fixes.

\begin{workedexample}
Imagine running ten populations, each starting with one copy of the focal allele. Suppose the outcomes, in order, are
\[
\text{lost, lost, fixed, lost, lost, lost, lost, fixed, lost, lost}.
\]
There are two fixation outcomes out of ten runs. Our estimate of the fixation probability is
\[
\frac{\text{number of fixations}}{\text{number of runs}}
=\frac{2}{10}=0.2.
\]
This is an invented small example to explain the calculation. It is not a result from the full study.
\end{workedexample}

We will use three symbols. The capital letter $R$ is the number of replicates. The lower-case letter $k$ is the number of those replicates that fix. The symbol $\widehat p$, read ``p hat'', is the estimated fixation probability:
\begin{equation}
\widehat p=\frac{k}{R}.
\label{eq:estimated-probability}
\end{equation}
The hat tells us that this number was estimated from a finite set of runs. We use $p$, without a hat, for the true fixation probability under the model. Usually a simulation study does not know $p$ in advance. Our model has an exact formula, which gives us a valuable way to check the simulation.

\section{What the repetition loop does}

The repetition loop is conceptually simple:
\begin{enumerate}
\item Set a fixation counter to zero.
\item Start a new population in the initial state.
\item Simulate until loss or fixation.
\item If fixation occurred, add one to the counter.
\item Repeat until $R$ populations have finished.
\item Divide the final counter by $R$.
\end{enumerate}

The \emph{outer} loop runs the independent experiments. Inside it, the simulation function contains an \emph{inner} loop that runs reproduction events for one experiment. Confusing these two loops is a common beginner mistake: twenty thousand events in one population are not twenty thousand independent populations.

Do not declare a long-running history to be a loss merely because it has not fixed yet. An unfinished history has an unknown eventual outcome. Our fixation experiments run until one of the two absorbing states is reached. We may use a shorter time window when drawing example histories, but that window is a display choice, not a rule for classifying fixation.

\section{Random seeds: repeatable does not mean identical trials}

A computer normally uses a \emph{pseudorandom-number generator}. This is a rule that produces a long sequence of numbers that behaves sufficiently like independent random draws for this simulation. A \emph{seed} sets its starting state.

If we use the same generator, seed, Julia version and code, we can reproduce the same sequence. This helps us check a calculation. It does not mean that successive replicates receive the same numbers: the generator moves forward each time a draw is requested.

Set the seed once before a batch of experiments. Do \emph{not} reset the generator to the same seed before every replicate. Doing that could make every replicate repeat the same history, creating an illusion of many experiments while adding no independent information.

Our study records a seed for each parameter combination. It also records the Julia version. Keeping these details with the results makes the calculation reproducible. Bit-for-bit random sequences are not promised across arbitrary changes of software version.

\begin{checkpoint}
\textbf{Question.} You reset the generator to seed 42, run a population, reset it to seed 42 again, and run the same population again. Have you obtained two useful independent trials?

\textbf{Worked answer.} No. With unchanged code and environment, both runs consume the same sequence of random numbers and produce the same history. Instead, create the generator once and pass that continuing generator into each replicate.
\end{checkpoint}

\section{Why estimates change from one batch to another}

Even a correct simulation produces slightly different estimates when we change its seed. That difference is \emph{Monte Carlo error}: uncertainty caused by using a finite number of random experiments. ``Monte Carlo'' here simply means estimating something by random sampling.

Monte Carlo error is different from \emph{model uncertainty}. Monte Carlo error asks whether we have run enough repetitions to estimate the probability of \emph{this model}. Model uncertainty asks whether the assumptions of this model adequately describe a biological situation. Increasing $R$ reduces the first problem. It does not remove the second.

We can describe Monte Carlo error mathematically. For replicate $j$, define a number $Y_j$ as follows:
\[
Y_j=
\begin{cases}
1,&\text{if replicate }j\text{ fixes},\\
0,&\text{if replicate }j\text{ is lost}.
\end{cases}
\]
The subscript $j$ is a label: $Y_1$ describes the first run, $Y_2$ the second, and so on. A random variable with possible values zero and one is called a \emph{Bernoulli random variable}.

The expected value, written $\mathbb E[Y_j]$, means its average over many repetitions. It equals each possible value multiplied by its probability:
\[
\mathbb E[Y_j]=1\times p+0\times(1-p)=p.
\]
Here $1-p$ is the probability of loss, because fixation and loss are the only eventual outcomes in our model.

The \emph{variance} measures squared fluctuations around an expected value. We write it as $\operatorname{Var}$. For a Bernoulli variable, squaring changes neither possible value: $0^2=0$ and $1^2=1$. Therefore
\[
\mathbb E[Y_j^2]=p,
\qquad
\operatorname{Var}(Y_j)
=\mathbb E[Y_j^2]-\bigl(\mathbb E[Y_j]\bigr)^2
=p-p^2=p(1-p).
\]

The total number of fixations is $K=Y_1+\cdots+Y_R$. We use capital $K$ when discussing this count before the experiment is run, and lower-case $k$ for its observed value. For independent replicates, variances add: the fluctuations of one run do not systematically push the others in the same direction. Thus
\[
\operatorname{Var}(K)=R\,p(1-p).
\]
Our estimate is $\widehat p=K/R$. Dividing a quantity by $R$ divides its squared fluctuations by $R^2$. Consequently,
\begin{align}
\operatorname{Var}(\widehat p)
&=\frac{\operatorname{Var}(K)}{R^2}\\
&=\frac{R\,p(1-p)}{R^2}\\
&=\frac{p(1-p)}{R}.
\label{eq:mc-variance}
\end{align}

The square root of this variance is the \emph{standard error}:
\begin{equation}
\operatorname{SE}(\widehat p)=\sqrt{\frac{p(1-p)}{R}}.
\label{eq:mc-se}
\end{equation}
It describes the typical size of the estimate's fluctuations across repeated batches. When $p$ is unknown, one common estimate replaces $p$ with $\widehat p$. When checking against our exact solution, we can calculate the standard error using the exact $p$.

Notice the square root. Multiplying the number of replicates by four approximately halves the standard error. Multiplying it by one hundred approximately divides the standard error by ten. More precision costs more computation.

\section{A rare outcome needs enough repetitions}

For neutral selection, one initial copy in a population of $N$ individuals fixes with probability $1/N$. In the largest population in our grid, $N=200$, so
\[
p=\frac{1}{200}=0.005.
\]
This is a half-percent fixation chance. With $R=20{,}000$ independent runs, the expected number of fixations is
\[
R p=20{,}000\times0.005=100.
\]
It need not be exactly 100. The standard error of the estimated probability is
\[
\sqrt{\frac{0.005\times0.995}{20{,}000}}
\approx0.000499.
\]
That is about 0.05 percentage points. A rough two-standard-error range around the true probability is about $0.005\pm0.001$. The absolute discrepancy is small, but $0.001/0.005=0.2$, or 20\% of the probability itself. For rare events, relative uncertainty can remain noticeable even when the absolute uncertainty looks tiny.

\begin{workedexample}
Suppose we use only 100 replicates when $p=0.005$. The expected fixation count is $100\times0.005=0.5$. Many batches will contain no fixations at all. Such a batch does not establish that fixation is impossible. It shows that 100 runs provide limited information about an event this rare.
\end{workedexample}

\section{Confidence intervals: show a range, not just a point}

A point estimate such as $\widehat p=0.1$ hides how much sampling information supports it. Ten fixations out of 100 runs and ten thousand fixations out of one hundred thousand runs both give 0.1, but the latter is much more precise.

A \emph{confidence interval} is a procedure that turns the observed count and sample size into a range. We use a \emph{Wilson interval}. It behaves more sensibly than the simplest symmetric interval when counts are small or the estimate is near zero or one.

The interval uses a number $z=1.96$, a conventional threshold for an approximate 95\% interval based on a bell-shaped reference distribution. You do not need to memorise 1.96 to understand what the code does. It is an input specifying the uncertainty level.

Define a denominator $D$, a centre $C$ and a half-width $H$:
\begin{align}
D&=1+\frac{z^2}{R},\\
C&=\frac{\widehat p+z^2/(2R)}{D},\\
H&=\frac{z}{D}\sqrt{\frac{\widehat p(1-\widehat p)}{R}
                           +\frac{z^2}{4R^2}}.
\end{align}
The lower endpoint is $C-H$ and the upper endpoint is $C+H$. The centre is generally a little different from $\widehat p$; that is intentional.

Where does this expression come from? Start with the comparison
\[
|\widehat p-p|\leq z\sqrt{\frac{p(1-p)}{R}}.
\]
The vertical bars mean ``absolute difference''. Temporarily treat $p$ as an unknown candidate probability. Squaring both sides and multiplying by $R$ gives
\[
R(\widehat p-p)^2\leq z^2p(1-p).
\]
Expand the square, bring everything to the left, and collect equal powers of $p$:
\[
(R+z^2)p^2-(2R\widehat p+z^2)p+R\widehat p^2\leq0.
\]
This is a quadratic inequality: it contains $p^2$, $p$ and a constant. Its two roots bound the acceptable values of $p$. Applying the quadratic formula gives exactly $C-H$ and $C+H$ above. The confidence level is approximate because the starting comparison uses an approximation to sampling fluctuations.

\begin{workedexample}
Take $R=100$ runs and $k=10$ fixations. First calculate $\widehat p=10/100=0.1$. With $z=1.96$, the Wilson calculation gives
\[
D=1.038416,\qquad C\approx0.114798,\qquad H\approx0.059569.
\]
The endpoints are approximately
\[
0.0552\quad\text{and}\quad0.1744.
\]
We would report an estimated fixation probability of 0.10, with a 95\% Wilson interval of about 0.055 to 0.174. These numbers illustrate the method; they do not come from the full parameter study.
\end{workedexample}

The 95\% statement describes the procedure across repeated batches: approximately 95\% of intervals produced in this way would contain the fixed model probability under suitable sampling conditions. It does not mean that, after calculating these endpoints, the fixed probability has a 95\% chance of lying inside them. Exact coverage can differ from 95\% because counts are discrete and the method is approximate.

If we observe zero fixations, the Wilson interval still has a positive upper endpoint. That is appropriate: observing no successes in a finite sample is compatible with a small positive probability.

\section{How this applies to our full study}

Our study has four population sizes and seven selection coefficients. That produces $4\times7=28$ parameter combinations. Each receives 20,000 independent replicates, for $28\times20{,}000=560{,}000$ complete population experiments.

For each combination, the output records the fixation count, point estimate, interval, exact prediction, seed and number of events used. The exact solution and the simulation are different ways to calculate the same model quantity. Their agreement is a check on our implementation.

Do not demand that all 28 nominal 95\% intervals contain their exact predictions. Even correct sampling can produce an occasional miss. Instead, examine whether discrepancies have the scale expected from Monte Carlo error, and use deterministic tests of the equations to check mathematical correctness. A systematic pattern of large discrepancies deserves investigation; one slightly missed interval is not by itself proof of a bug.

The full study and the quick demonstration save their outputs separately. A quick run is useful for learning and checking that the code starts correctly. Its smaller sample size must not silently replace the full study in the guide's results.

\begin{checkpoint}
\textbf{Question 1.} There are 37 fixations in 1,000 runs. What is the point estimate?

\textbf{Answer.} Divide fixations by runs: $\widehat p=37/1000=0.037$, or 3.7\%.

\medskip
\textbf{Question 2.} You want to halve the standard error without changing the model. Roughly how should you change the number of runs?

\textbf{Answer.} Use four times as many. In equation~\eqref{eq:mc-se}, replacing $R$ with $4R$ multiplies the denominator inside the square root by four, dividing the standard error by two.

\medskip
\textbf{Question 3.} Would one million repetitions prove that this model accurately describes a particular animal population?

\textbf{Answer.} No. More repetitions estimate the model's answer more precisely. They do not test whether its assumptions match that animal population. That requires biological evidence and possibly a different model.
\end{checkpoint}
